
\subsection{Segmentation improvements}
\label{sec:gap-obtaining}

Automated segmentation matches expert annotators on overlap and does not match them on topology.
The benchmark published with the dataset used here scored eight methods against two analysts on 160
CCTA scans: the best method, CAS-Net, reached a Dice similarity coefficient (DSC) of 91.2 against the
analysts' 92.8, but a Betti error (the number of connected components and loops by which the
predicted tree differs from the true one) of 1.9 against their 0.4 \cite{bransby2026imagecasx,
dong2023casnet}. The overlap gap is 1.6 points; the topology gap is nearly five-fold.

Overlap cannot see that difference, and optimizing for topology directly does not close it either.
clDice scores agreement along the vessel's skeleton rather than its volume and provably preserves
topology up to homotopy equivalence \cite{shit2021cldice}, so two masks can share a DSC and differ in
tree structure. Yet in the same benchmark, nnU-Net trained with an added clDice term gained 0.2 DSC
over plain nnU-Net and returned the table's worst Betti error, 8.0 \cite{isensee2021nnunet}, and the
two leading methods, CAS-Net and ADE-HTL, separate on DSC and Betti error ($p<0.001$) but tie on
clDice and both centerline distance measures \cite{zhang2024adehtl}. The benchmark's own qualitative
reading is that vessel breaks appear in all model predictions despite high DSC and clDice
\cite{bransby2026imagecasx}.

A break costs little overlap and can cost a whole feature: losing a distal branch moves DSC by a
fraction of a point, removes a branch from a branch count, and in the posterior vessels can change
which artery is called dominant.

This thesis does not train or fine-tune a segmentation network. CAS-Net's published weights are used
as delivered \cite{bransby2026imagecasx}; the effort goes into recovering topological correctness
afterward, as post-processing on the predicted mask or centerline, not as a training objective. Qiu
et al.\ take the same view: a centerline-enhanced loss, a distance-and-direction reconnection step and
implicit reconstruction of missing vessels together reach a DSC of 88.5 and 85.1 on ASOCA and PDSCA
\cite{qiu2025topology} -- a different cohort, so the numbers do not transfer directly, but evidence
that a post-hoc repair stage is still needed after a well-trained network, which is the stage this
thesis targets.